% Part: proof-theory
% Chapter: proof-search
% Section: tableaux

\documentclass[../../../include/open-logic-section]{subfiles}

\begin{document}

\olfileid[pt]{pt}{ps}{tab}

\olsection{Tableaux}

Os sistemas de tableaux são sistemas de !!{proof} estreitamente
relacionados aos cálculos de sequentes e particularmente adequados à
busca de provas, tanto manualmente quanto em implementações de software.
Em vez de construir árvores de sequentes, nos tableaux !!a{proof} é uma
árvore de !!{formula}s. No entanto, ao contrário da dedução natural, as
regras refletem as do cálculo de sequentes. Uma !!{proof} por tableau é,
em essência, uma busca explícita por um modelo. Por isso, os tableaux às
vezes são chamados de tableaux \emph{semânticos} ou também de
\emph{árvores de verdade}.

Um tableau \emph{com sinais} é uma árvore de !!{formula}s com sinais,
que podem ter a forma $\sFmla{\True}{!A}$ ou $\sFmla{\False}{!A}$. As
árvores crescem para baixo. Elas começam com !!{formula}s que especificam
o que queremos !!{prove}. Por exemplo, para !!{prove} um sequente $!A,
!B \Sequent !C, !D$, começaríamos com $\sFmla{\True}{!A},
\sFmla{\True}{!B}, \sFmla{\False}{!C}, \sFmla{\False}{!D}$.
!!^a{structure} que torna $!A$ e~$!B$ verdadeiras e $!C$ e~$!D$ falsas
é !!a{structure} que mostra que $!A, !B \Sequent !C, !D$ não é
válido. 

Um nó folha (o mais baixo) do tableau pode ser estendido por regras de
extensão de ramo, que acrescentam outros nós ao mesmo ramo ou acrescentam
dois nós irmãos abaixo dele, dividindo o ramo em dois novos ramos. As
regras podem ser aplicadas a qualquer !!{formula} com sinal que esteja
acima no mesmo ramo. As regras de extensão de ramo do sistema clássico
de tableaux \Log{Tc} são apresentadas em \olref{tab:Tc}.

\subfile{rules-Tc}

Como se pode ver, as regras são simplesmente as de \Log{G3c} de cabeça
para baixo, e os sinais $\True$ e $\False$ indicam, respectivamente, se
!!a{formula} é principal no lado esquerdo ou direito de um sequente.

Um ramo de um tableau é \emph{fechado} se contém tanto
$\sFmla{\True}{!A}$ quanto $\sFmla{\False}{!A}$. Se todos os ramos são
fechados, dizemos que o tableau inteiro está fechado. Por exemplo, este
é um tableau de cláusulas para $!A \lor (!B \land !C) \Sequent (!A \lor
!B) \land (!A \lor !C)$, isto é, que começa com
$\sFmla{\True}{!A \lor (!B \land !C)}$ e
$\sFmla{\False}{(!A \lor !B) \land (!A \lor !C)}$:
\begin{maxwidth}
  \begin{tableau}{}
  [\sFmla{\True}{\formula{A} \lor (\formula{B} \land \formula{C})},
    just=\TAss, checked
    [\sFmla{\False}{(\formula{A} \lor \formula{B}) \land
        (\formula{A} \lor \formula{C})}, just=\TAss, checked
      [\sFmla{\True}{\formula{A}}, just={\TRule{\True}{\lor}[1]}
        [\sFmla{\False}{\formula{A} \lor \formula{B}},
          just={\TRule{\False}{\land}[2]}, checked
          [\sFmla{\False}{\formula{A}},
            just={\TRule{\False}{\lor}[4]}
            [\sFmla{\False}{\formula{B}},
              just={\TRule{\False}{\lor}[4]}, close]
          ]
        ]
        [\sFmla{\False}{\formula{A} \lor \formula{C}},
          just={\TRule{\False}{\land}[2]}, checked
          [\sFmla{\False}{\formula{A}},
            just={\TRule{\False}{\lor}[4]}
            [\sFmla{\False}{\formula{C}},
              just={\TRule{\False}{\lor}[4]}, close]
          ]
        ]
      ]
      [\sFmla{\True}{\formula{B} \land \formula{C}},
        just={\TRule{\True}{\lor}[1]}, checked
        [\sFmla{\False}{\formula{A} \lor \formula{B}},
          just={\TRule{\False}{\land}[2]}, checked
          [\sFmla{\True}{\formula{B}},
            just={\TRule{\True}{\land}[3]}, move by=2
            [\sFmla{\True}{\formula{C}},
              just={\TRule{\True}{\land}[3]}
              [\sFmla{\False}{\formula{A}},
                just={\TRule{\False}{\lor}[4]}
                [\sFmla{\False}{\formula{B}},
                  just={\TRule{\False}{\lor}[4]},close
                ]
              ]
            ]
          ]
        ]
        [\sFmla{\False}{\formula{A} \lor \formula{C}},
          just={\TRule{\False}{\land}[2]}, checked
          [\sFmla{\True}{\formula{B}},
            just={\TRule{\True}{\land}[3]}, move by=2
              [\sFmla{\True}{\formula{C}},
                just={\TRule{\True}{\land}[3]}
                [\sFmla{\False}{\formula{A}},
                  just={\TRule{\False}{\lor}[4]}
                  [\sFmla{\False}{\formula{C}},
                  just={\TRule{\False}{\lor}[4]},close
                ]
              ]
            ]
          ]
        ]
      ]
    ]
  ]
\end{tableau}
\end{maxwidth}

As marcas de verificação indicam que a regra correspondente foi aplicada
a !!a{formula} em todos os ramos abaixo dela. O símbolo~$\otimes$ indica
que um ramo está fechado. Os números de linha após as regras indicam a
quais !!{formula}s com sinais a regra é aplicada (a saber, à !!{formula}
no mesmo ramo e na linha indicada).

A busca de provas é simples em tableaux. Geramos um tableau em etapas.
Na etapa~$0$, apenas escrevemos no topo as !!{formula}s iniciais. Na
etapa~$n+1$, fazemos o seguinte para cada nó folha: encontramos a primeira
fórmula com sinal (a mais alta) que ainda não foi marcada; aplicamos a
regra de extensão de ramo a essa fórmula; depois que a regra tiver sido
aplicada em todo ramo que contém a fórmula com sinal em questão, nós a
marcamos. (O exemplo acima é gerado com esse algoritmo.)

As !!{formula}s com sinais da forma $\sFmla{\True}{\lforall}$ e
$\sFmla{\False}{\lexists}$ nunca são marcadas e, portanto, precisam de
tratamento especial. Se estiverem presentes, precisamos garantir que as
regras $\TRule{\True}{\forall}$ e $\TRule{\False}{\lexists}$ sejam
finalmente aplicadas a todos os termos pertinentes. Podemos garantir isso,
por exemplo, aplicando-as em cada etapa a todas essas fórmulas em todos os
ramos (depois de tratar a !!{formula} com sinal mais alta que não tenha
essa forma). O termo $t$ usado é o primeiro termo~$t$ que pode ser
construído a partir das !!{constant}s que já ocorrem no ramo (ou $\Obj
c_0$, se nenhuma !!{constant} ocorre) e das !!{function}s que ocorrem nas
!!{formula}s iniciais, e tal que a !!{formula} com sinal correspondente
$\sFmla{\True}{!B(t)}$ ou~$\sFmla{\False}{!B(t)}$ ainda não ocorra no
ramo.

Se essa busca de provas não produz um tableau fechado, ele contém um
ramo aberto finito no qual todas as !!{formula}s não atômicas estão
marcadas, ou então a busca produz uma árvore infinita com ramificação
finita, que necessariamente contém um ramo infinito. Assim como em
\olref[cpl]{sec}, podemos definir !!a{structure}~$\Struct{M}$ tal que,
se $\sFmla{\True}{!A}$ ocorre no ramo, então $\Sat{M}{!A}$, e, se
$\sFmla{\False}{!A}$ ocorre, então $\Sat/{M}{!A}$. (Tomamos $\Theta =
\Setabs{!A}{\sFmla{\True}{!A} \text{ no ramo}}$ e $\Xi =
\Setabs{!A}{\sFmla{\False}{!A} \text{ no ramo}}$.)

Os tableaux gerados por esse método de busca de provas podem ser
traduzidos facilmente em provas de \Log{G3c}, atribuindo-se um sequente
a cada nó. Se as fórmulas iniciais correspondem a um sequente $\Gamma
\Sequent \Delta$, rotulamos cada uma delas com $\Gamma \Sequent
\Delta$. Se um ramo é estendido pela aplicação de uma regra de extensão
a uma !!{formula} com sinal $\sFmla{\True}{!A}$, a folha recebe o rótulo
$!A, \Pi \Sequent \Lambda$; se a regra é aplicada a uma !!{formula} com
sinal $\sFmla{\False}{!A}$, a folha recebe o rótulo $\Pi \Sequent
\Lambda, !A$. Quando o tableau aplica uma regra
\TRule{\True}{\star}, aplicamos ao sequente a regra \LeftR{\star} (com
$!A$ principal); quando ele aplica uma regra \TRule{\False}{\star},
aplicamos a regra \RightR{\star}. As premissas da regra são os sequentes
que rotulam os nós correspondentes acrescentados ao tableau. Quando um
tableau se fecha com $\sFmla{\True}{!A}$ e $\sFmla{\False}{!A}$, o
rótulo do nó folha final será um sequente da forma $!A, \Pi \Sequent
\Lambda, !A$. Alguns nós sucessivos no tableau receberão o mesmo
sequente; eliminamos as duplicatas. Por exemplo, o tableau acima se
traduz na !!{proof}:

\begin{maxwidth}
$
\Axiom$!A \fCenter !A, !B$
\RightLabel{\RightR{\lor}}
\UnaryInf$!A \fCenter !A \lor !B$
\Axiom$!A \fCenter !A, !C$
\RightLabel{\RightR{\lor}}
\UnaryInf$!A \fCenter !A \lor !C$
\RightLabel{\RightR{\land}}
\BinaryInf$!A \fCenter (!A \lor !B) \land (!A \lor !C)$
\Axiom$!B, !C \fCenter !A, !B$
\RightLabel{\RightR{\lor}}
\UnaryInf$!B, !C \fCenter !A \lor !B$
\RightLabel{\LeftR{\land}}
\UnaryInf$!B \land !C \fCenter !A \lor !B$
\Axiom$!B, !C \fCenter !A, !C$
\RightLabel{\RightR{\lor}}
\UnaryInf$!B, !C \fCenter !A \lor !C$
\RightLabel{\LeftR{\land}}
\UnaryInf$!B \land !C \fCenter !A \lor !C$
\RightLabel{\RightR{\land}}
\BinaryInf$!B \land !C \fCenter (!A \lor !B) \land
(!A \lor !C)$
\RightLabel{\LeftR{\lor}}
\BinaryInf$!A \lor (!B \land !C) \fCenter (!A \lor !B) \land
(!A \lor !C)$
\DisplayProof$
\end{maxwidth}

\end{document}
